\exercisesection{}

\begin{enumerate}
\def\labelenumi{\arabic{enumi})}
\tightlist
\item
  Let G be a compact topological group and let B be the family of characters of irreducible unitary representations of G.
\end{enumerate}

\begin{enumerate}
\def\labelenumi{\alph{enumi})}
\item
  The pair \((\mathrm{R}(\mathrm{G}),\mathcal{B})\) is a natural unital basal algebra (cf.~exercise 6 of III, p.~130).
\item
  Assume that G is finite. The algebra R(G) is then of finite rank; compute \(p_f(\pi)\) for every \(\pi \in \widehat{G}\), and identify the representation of G corresponding to the element r₀ defined in question j) of loc. cit.
\end{enumerate}

\begin{enumerate}
\def\labelenumi{\arabic{enumi})}
\setcounter{enumi}{1}
\tightlist
\item
  Let \(n \geqslant 1\) and \(p \in [1, +\infty]\) . We equip \(\mathbf{R}^n\) with the norm
\end{enumerate}

\[
\| x \| = \left(\sum_ {i = 1} ^ {n} | x _ {i} | ^ {p}\right) ^ {1 / p} \text {   if   } p \neq + \infty ,
\]

and \(\|x\| = \sup_{i}|x_{i}|\) if \(p = +\infty\) . Let G be the group of isometries of the metric space \(R^{n}\) equipped with this norm.

\begin{enumerate}
\def\labelenumi{\alph{enumi})}
\item
  Let \(\sigma\) be a permutation of \(\{1,\ldots,n\}\) and \(\varepsilon\in\{-1,1\}^{n}\) . Let \((e_{i})\) be the canonical basis of \(\mathbf{R}^{n}\) . The unique endomorphism of \(\mathbf{R}^{n}\) such that \(e_{i}\mapsto\varepsilon_{i}e_{\sigma(i)}\) belongs to G. The set of isometries of this form is a subgroup W of G; it is an extension of \(\mathfrak{S}_{n}\) by \((\mathbf{Z}/2\mathbf{Z})^{n}\) .
\item
  The group G is compact.
\item
  The action of the group G on \(R^{n}\) is irreducible; the group G is contained in the orthogonal group of the quadratic form \(\sum x_{i}^{2}\) .
\item
  If \(p \neq 2\) , then G = W.
\end{enumerate}

\begin{enumerate}
\def\labelenumi{\arabic{enumi})}
\setcounter{enumi}{2}
\tightlist
\item
  We recall that a group \(\Gamma\) is said to be residually finite if the intersection of the finite-index subgroups of \(\Gamma\) is reduced to \(\{e\}\) (A, I, p.~129, exerc. 5, b)).
\end{enumerate}

For a group \(\Gamma\) , we denote by \(\widehat{\Gamma}\) the projective limit \(\widehat{\Gamma} = \varprojlim \Gamma / H\) , where \(H\) ranges over the set of finite-index normal subgroups of \(\Gamma\) ordered by inclusion. We say that \(\widehat{\Gamma}\) is the profinite completion of \(\Gamma\) .

\begin{enumerate}
\def\labelenumi{\alph{enumi})}
\item
  A group \(\Gamma\) is residually finite if and only if, for every \(g \in \Gamma\) different from the identity element, there exists a finite group F and a homomorphism \(\varphi: \Gamma \to F\) such that \(\varphi(g) \neq e\) .
\item
  Show that \(\Gamma\) is residually finite if and only if the canonical homomorphism from \(\Gamma\) into \(\widehat{\Gamma}\) is injective.
\item
  Let K be a field. Every finitely generated subgroup (cf.~A, I, p.~145, exerc. 16) of \(\mathbf{GL}(n,\mathrm{K})\) is residually finite ("Mal'cev's theorem"). (Show that there exists an integer \(m\geqslant 0\) , a ring A equal to \(\mathbf{Z}\) or to a finite field and an injective homomorphism of such a subgroup into \(\mathbf{GL}(m,\mathrm{R})\) , where R is of the form \(\mathrm{A}[\mathrm{X}_1,\dots ,\mathrm{X}_r,1 / f]\) .)
\item
  Let \(\Gamma\) be a finitely generated discrete group. Suppose that there exists an injective homomorphism from \(\Gamma\) into a compact topological group. For every \(\gamma \in \Gamma\) different from the identity element, there exists a finite-dimensional unitary irreducible representation of \(\Gamma\) such that the image of \(\gamma\) by \(\pi\) is not the identity map.
\item
  A finitely generated group is residually finite if and only if it is isomorphic to a subgroup of a compact group.
\end{enumerate}

¶ 4) Let G be a group and let A and B be subgroups of G. Let \(\varphi\colon A\to B\) be an isomorphism.

Let \(r'\) be a family in the free group F(G) such that G is isomorphic to the group F(G, \(r'\) ) (A, I, p.~87, n°6). One calls the HNN extension of G according to \(\varphi\) the group \(G^{\varphi} = F(X, r)\) where X is the sum of the sets G, A, B and an element t, and where r is the sum family of \(r'\) and of the relators \(t^{-1}at = \varphi(a)\) for all \(a \in A\) .

We denote by \(j\colon G\to G^{\varphi}\) the unique homomorphism such that \(j(g)\) is the element represented by \(g\) in \(G^{\varphi}\) .

\begin{enumerate}
\def\labelenumi{\alph{enumi})}
\tightlist
\item
  Let \(k \geqslant 1\) an integer and let \((g_{0}, \varepsilon_{1}, g_{1}, \ldots, \varepsilon_{k}, g_{k})\) be a finite family such that \(g_{i} \in G\) and \(\varepsilon_{i} \in \{-1, 1\}\) for every i. Suppose that there does not exist an integer j such that \(1 \leqslant j \leqslant k - 1\) and
\end{enumerate}

\[
(\varepsilon_ {j}, \varepsilon_ {j + 1}) = (- 1, 1), \quad g _ {i} \in A
\]

or

\[
(\varepsilon_ {j}, \varepsilon_ {j + 1}) = (1, - 1), \quad g _ {i} \in \mathrm{B}.
\]

Then the element

\[
j (g _ {0}) t ^ {\varepsilon_ {1}} j (g _ {1}) \dots t ^ {\varepsilon_ {k}} j (g _ {k})
\]

of G \(^{φ}\) is not the identity element ("Britton's Lemma".)

\begin{enumerate}
\def\labelenumi{\alph{enumi})}
\setcounter{enumi}{1}
\item
  The homomorphism j is injective.
\item
  Let \(a\) and \(b\) be distinct symbols. Let G be the group \(\mathrm{F}(\{a,b\}, a^{-1}b^{2}ab^{-3})\) . We denote by \(b_{1} = a^{-1}ba\) . Let H be a normal subgroup of finite index of G and \(\pi: \mathrm{G} \to \mathrm{G}/\mathrm{H}\) the canonical projection. One has \([b_{1}, b] \in \mathrm{Ker}(\pi)\) .
\item
  The group G admits no faithful finite-dimensional representation. (Show that G is not residually finite.)
\end{enumerate}

\begin{enumerate}
\def\labelenumi{\arabic{enumi})}
\setcounter{enumi}{4}
\tightlist
\item
  Let G be a compact group and \(\varrho\) a faithful finite-dimensional unitary representation of G. Let \(\pi \in \widehat{G}\) .
\end{enumerate}

\begin{enumerate}
\def\labelenumi{\alph{enumi})}
\tightlist
\item
  There exist integers a and b such that \(\pi\) is isomorphic to a subrepresentation of \(\varrho^{\otimes a} \otimes \overline{\varrho}^{\otimes b}\) . (One may use, for example, the Weierstrass--Stone theorem.)
\end{enumerate}

We denote by \(w(\pi)\) the minimum of \(a + b\) when \((a, b)\) ranges over the set of integers satisfying this property.

\begin{enumerate}
\def\labelenumi{\alph{enumi})}
\setcounter{enumi}{1}
\tightlist
\item
  Let \(\Lambda\) be the weight lattice of \(G^{0}\) and let q be a norm on \(\Lambda\otimes \mathbf{R}\) . There exist real numbers C and D such that for every representation \(\pi\in\widehat{G}\) whose restriction to \(G^{0}\) has highest weights \((\lambda_{1},\ldots,\lambda_{k})\) , we have
\end{enumerate}

\[
w (\pi) \leqslant \mathrm{C} \sup (q (\lambda_ {1}), \dots , q (\lambda_ {k})) + \mathrm{D}.
\]

(Consider first the case of a connected semisimple group and apply a) to a suitably chosen finite family of representations of G).

\begin{enumerate}
\def\labelenumi{\arabic{enumi})}
\setcounter{enumi}{5}
\tightlist
\item
  Let G be a compact group and let \(\varrho\) be an irreducible representation of G. There exists a central function \(f \in \Theta(G^{\sharp})\) such that
\end{enumerate}

\[
f = \sum_ {\pi \in \widehat {\mathrm{G}}} c (\pi) \chi_ {\pi}
\]

where the integers \(c(\pi)\) satisfy

\begin{enumerate}
\def\labelenumi{(\roman{enumi})}
\item
  \(c(\pi) \geqslant 0\) for all \(\pi \in \widehat{G}\) ;
\item
  \(c(1) \leqslant c(\varrho)\) , and the inequality is strict except possibly if \(\varrho\) is of dimension 1 and satisfies \(\varrho^{2} = 1\) ;
\item
  \(\mathcal{R}(f) \geqslant 0.\)
\end{enumerate}

¶ 7) For every integer \(n \geqslant 1\) , we denote by \(\mu_n\) the positive measure on \(\mathbf{C}\) image of the normalized Haar measure of the group \(\mathbf{U}(\mathbf{C}^n) \subset \mathrm{GL}(\mathbf{C}^n)\) under the trace map \(g \mapsto \mathrm{Tr}(g)\) ; it is a compactly supported measure of total mass 1.

\begin{enumerate}
\def\labelenumi{\alph{enumi})}
\tightlist
\item
  Let \(a\) and \(b\) be natural integers such that \(a + b \leqslant n\) . We have
\end{enumerate}

\[
\int_ {\mathbf {C}} z ^ {a} \overline {{z}} ^ {b} d \mu_ {n} (z) = \left\{ \begin{array}{l l} 0 & \text{if } a \neq b \\ a! & \text{if } a = b. \end{array} \right.
\]

\begin{enumerate}
\def\labelenumi{\alph{enumi})}
\setcounter{enumi}{1}
\tightlist
\item
  The sequence of measures \((\mu_{n})_{n\geqslant1}\) converges narrowly to the Gaussian measure on the real vector space C whose covariance is \(\mathrm{Q}(z)=|z|^{2}\) .
\end{enumerate}

\begin{enumerate}
\def\labelenumi{\arabic{enumi})}
\setcounter{enumi}{7}
\item
  Let G be a compact group and let \(\varrho_{1}\) and \(\varrho_{2}\) be irreducible unitary representations of G. Let \(n \geqslant 2\) be an integer. If \(S^{n}\varrho_{1}\) is isomorphic to \(S^{n}\varrho_{2}\) , then there exists a subgroup H of finite index of G such that the restriction of \(\varrho_{1}\) to H is isomorphic to the restriction of \(\varrho_{2}\) to H.
\item
  Let \(k\) be a finite field with 5 elements and \(\mathrm{G} = \mathbf{S}\mathbf{L}(2,k)\) .
\end{enumerate}

\begin{enumerate}
\def\labelenumi{\alph{enumi})}
\tightlist
\item
  The conjugacy classes of the group G are represented by e, -e and the seven elements
\end{enumerate}

\[
g _ {3} = \left( \begin{array}{c c} 1 & 1 \\ 0 & 1 \end{array} \right), \qquad g _ {4} = \left( \begin{array}{c c} 1 & 2 \\ 0 & 1 \end{array} \right), \qquad g _ {5} = \left( \begin{array}{c c} - 1 & 1 \\ 0 & - 1 \end{array} \right), \qquad g _ {6} = \left( \begin{array}{c c} - 1 & 2 \\ 0 & - 1 \end{array} \right)
\]

\[
g _ {7} = \left( \begin{array}{c c} 2 & 0 \\ 0 & - 2 \end{array} \right), \qquad g _ {8} = \left( \begin{array}{c c} 0 & 1 \\ - 1 & 1 \end{array} \right), \qquad g _ {9} = \left( \begin{array}{c c} 0 & 1 \\ - 1 & - 1 \end{array} \right).
\]

Determine the cardinality of each conjugacy class.

\begin{enumerate}
\def\labelenumi{\alph{enumi})}
\setcounter{enumi}{1}
\tightlist
\item
  Let \(\alpha = \frac{1}{2}(1 + \sqrt{5})\) . There exists a faithful irreducible unitary representation \(\pi\) of dimension 2 of G whose character \(\chi\) satisfies \(\chi(-e) = -2\) and
\end{enumerate}

\[
\begin{array}{c} \chi (g _ {4}) = \alpha , \quad \chi (g _ {6}) = - \alpha , \quad \chi (g _ {3}) = - \alpha - 1, \quad \chi (g _ {5}) = \alpha + 1 \\ \chi (g _ {7}) = 0, \quad \chi (g _ {8}) = 1, \quad \chi (g _ {9}) = - 1. \end{array}
\]

(One may consider the restriction to G of the irreducible representations of dimension 4 of \(\mathbf{GL}(2,k)\), cf.~A, VIII, p.~422, exercise 27.)

\begin{enumerate}
\def\labelenumi{\alph{enumi})}
\setcounter{enumi}{2}
\tightlist
\item
  One has \(\mathrm{M}_4(\pi) = 2\) .
\end{enumerate}

\begin{enumerate}
\def\labelenumi{\arabic{enumi})}
\setcounter{enumi}{9}
\tightlist
\item
  \begin{enumerate}
  \def\labelenumii{\alph{enumii})}
  \tightlist
  \item
    There exists an irreducible representation \(\pi\) of orthogonal type of the alternating group \(\mathfrak{A}_5\) such that \(\mathrm{M}_4(\pi) = 3\) .
  \end{enumerate}
\end{enumerate}

\begin{enumerate}
\def\labelenumi{\alph{enumi})}
\setcounter{enumi}{1}
\tightlist
\item
  Let R be a root system of type \(E_{6}\) (resp. of type \(E_{7}\), of type \(E_{8}\) ) and W the Weyl group of R. Let \(\pi\) be the geometric representation of W (LIE, V, § 4, no. 3). It is irreducible, of orthogonal type, and satisfies \(M_{4}(\pi)=3\) .
\end{enumerate}

(One may use software such as GAP or MAGMA to determine the list of conjugacy classes of W and their images under \(\pi\) .)

\begin{enumerate}
\def\labelenumi{\arabic{enumi})}
\setcounter{enumi}{10}
\item
  Let G be a compact group and \(\pi\) a finite-dimensional unitary representation of G in a Hilbert space E. If \(\mathrm{M}_{4}(\pi) \leqslant 5\) then the representation \(\pi\) is irreducible.
\item
  Let G be a connected compact semisimple Lie group whose complexified Lie algebra is isomorphic to the simple Lie algebra \(e_{6}\) .
\end{enumerate}

\begin{enumerate}
\def\labelenumi{\alph{enumi})}
\item
  There exists a faithful irreducible representation \(\varrho\) of G of dimension 27; such a representation is of complex type.
\item
  The fourth moment of the representation \(\varrho\) is equal to 3.
\end{enumerate}

\begin{enumerate}
\def\labelenumi{\arabic{enumi})}
\setcounter{enumi}{12}
\tightlist
\item
  Let G be a connected compact semisimple Lie group and let g be its Lie algebra. Suppose there exists a faithful irreducible representation \(\varrho\) of dimension 27 of G. We assume that the fourth moment of \(\varrho\) is equal to 3 and that the representation \(\varrho\) is of complex type.
\end{enumerate}

\begin{enumerate}
\def\labelenumi{\alph{enumi})}
\item
  The Lie algebra g is simple.
\item
  The Lie algebra g is isomorphic to \(\mathfrak{e}_{6}\) .
\end{enumerate}

\begin{enumerate}
\def\labelenumi{\arabic{enumi})}
\setcounter{enumi}{13}
\item
  Let G be a connected compact semisimple Lie group. Suppose there exists a faithful irreducible representation \(\varrho\) of dimension 27 of G. The complexified Lie algebra of G is isomorphic to \(\mathfrak{e}_6\) if and only if there exists an irreducible subrepresentation of dimension 78 in the representation \(\operatorname{End}(\varrho)\) .
\item
  Let \(q\) be a power of an odd prime number and F a finite field with \(q\) elements. Consider the Heisenberg group H \(_3\) (F) of matrices of size 3 with coefficients in F that are upper triangular and unipotent. For \((a, b, c) \in \mathrm{F}^3\) , we denote
\end{enumerate}

\[
\sigma (a, b, c) = \left( \begin{array}{c c c} 1 & a & c \\ 0 & 1 & b \\ 0 & 0 & 1 \end{array} \right) \in \mathrm{H} _ {3} (\mathrm{F}).
\]

The map \(\sigma\) is therefore a bijection.

\begin{enumerate}
\def\labelenumi{\alph{enumi})}
\item
  The center C of \(\mathrm{H}_{3}(\mathrm{F})\) is identified with F by the group homomorphism sending \(\sigma(a,b,c)\) to c.
\item
  There exist \(q^{2}\) one-dimensional unitary characters of \(\mathrm{H}_{3}(\mathrm{F})\) ; these characters are trivial on C.
\item
  Let X be the subgroup of index q of F consisting of the elements \(\sigma(0,b,c)\) for \((b,c)\in\mathrm{F}^{2}\) . The subgroup X is identified with \(F^{2}\) . Let \(\psi_{1}\) and \(\psi_{2}\) be unitary characters of F, and set \(\psi(\sigma(0,b,c))=\psi_{1}(b)\psi_{2}(c)\) . The unitary representation \(\varrho\) of \(\mathrm{H}_{3}(\mathrm{F})\) induced by the character \(\psi\) of X is irreducible if and only if \(\psi_{2}\) is nontrivial; it depends only on \(\psi_{2}\) up to isomorphism, and its central character is equal to \(\psi_{2}\) . The support of the character of \(\varrho\) is equal to C.
\item
  The irreducible representations constructed in the two preceding questions are all the irreducible unitary representations of \(\mathrm{H}_{3}(\mathrm{F})\) .
\end{enumerate}

\begin{enumerate}
\def\labelenumi{\alph{enumi})}
\setcounter{enumi}{4}
\tightlist
\item
  Let \(\varrho\) be one of the irreducible unitary representations of dimension q of \(\mathrm{H}_{3}(\mathrm{F})\) . The representation \(\varrho\) is faithful and satisfies \(\mathrm{M}_{4}(\varrho)=q^{2}\) .
\end{enumerate}

\begin{enumerate}
\def\labelenumi{\arabic{enumi})}
\setcounter{enumi}{15}
\tightlist
\item
  Let \(r \geqslant 1\) be an integer and let G be a connected compact subgroup of \(\mathbf{U}(\mathbf{C}^{r})\) acting irreducibly on \(C^{r}\) . We denote by D the subgroup of diagonal matrices in \(\mathbf{U}(\mathbf{C}^{r})\) .
\end{enumerate}

\begin{enumerate}
\def\labelenumi{\alph{enumi})}
\tightlist
\item
  Let A be a subgroup of D contained in the normalizer of G in \(\mathbf{U}(\mathbf{C}^{r})\) . Let S be the subset of \(\widehat{A}\) formed by the diagonal coefficients \(a = (a_{i,j}) \mapsto a_{i,i}\) , where \(1 \leqslant i \leqslant r\) . We assume that S contains r elements and that the equation
\end{enumerate}

\[
\chi_ {1} \chi_ {2} ^ {- 1} = \chi_ {3} \chi_ {4} ^ {- 1}
\]

with \(\chi_{i} \in S\) admits only the solutions of the form \((\chi_{1}, \chi_{1}, \chi_{3}, \chi_{3})\) and \((\chi_{1}, \chi_{2}, \chi_{1}, \chi_{2})\) (one then says that \(S\) is a Sidon set).

If G is semisimple, prove that \(\mathbf{G} = \mathbf{S}\mathbf{U}(\mathbf{C}^{r})\) (Consider the representation of A on \(\operatorname{End}(\mathbf{C}^{r})\) , and deduce from the hypothesis that the complexified Lie algebra of G is stable under conjugation by D.)

\begin{enumerate}
\def\labelenumi{\alph{enumi})}
\setcounter{enumi}{1}
\tightlist
\item
  The group G contains \(\mathbf{SU}(\mathbf{C}^{r})\) if and only if there exists an element \(g \in G\) such that the eigenvalues of g are distinct and form a Sidon set in \(C^{*}\) .
\end{enumerate}

\begin{enumerate}
\def\labelenumi{\arabic{enumi})}
\setcounter{enumi}{16}
\tightlist
\item
  Let G be a compact group and H a closed subgroup of G. Let \(\varrho\) be a unitary representation of H in a finite-dimensional Hilbert space E.
\end{enumerate}

\begin{enumerate}
\def\labelenumi{\alph{enumi})}
\tightlist
\item
  Let \(\pi \in \widehat{\mathbf{G}}\) . The space \(\mathrm{Hom}_{\mathrm{G}}(\pi, \mathrm{Ind}_{\mathrm{H}}^{\mathrm{G}}(\varrho))\) is of finite dimension.

\end{enumerate}

\begin{enumerate}
\def\labelenumi{\alph{enumi})}
\setcounter{enumi}{1}
\item
  For every \(u \in \mathrm{Hom}_{\mathrm{G}}(\pi, \mathrm{Ind}_{\mathrm{H}}^{\mathrm{G}}(\varrho))\) , the image of \(u\) is contained in the image of \(\mathcal{C}(\mathrm{G};\mathrm{E})\) in \(\mathrm{L}^2 (\mathrm{G};\mathrm{E})\) .
\item
  There exists a linear map \(\Phi\) from the space \(\mathrm{Hom}_{\mathrm{G}}(\pi, \mathrm{Ind}_{\mathrm{H}}^{\mathrm{G}}(\varrho))\) into the space \(\mathrm{Hom}_{\mathrm{H}}(\mathrm{Res}_{\mathrm{H}}^{\mathrm{G}}(\pi), \varrho)\) such that, for every \(u \in \mathrm{Hom}_{\mathrm{G}}(\pi, \mathrm{Ind}_{\mathrm{H}}^{\mathrm{G}}(\varrho))\) and every \(x \in \mathrm{E}_{\pi}\) , \(\Phi(u)(x)\) is the image at \(e\) of the continuous function representing \(u(x)\) .
\item
  The linear map \(\Phi\) is an isomorphism.
\item
  Suppose that \(\varrho\) is irreducible. The multiplicity of \(\pi\) in \(\mathrm{Ind}_{\mathrm{H}}^{\mathrm{G}}(\varrho)\) is equal to the multiplicity of \(\varrho\) in \(\mathrm{Res}_{\mathrm{H}}^{\mathrm{G}}(\pi)\) (Frobenius reciprocity formula, cf. A, VIII, p. 392).
\end{enumerate}

\begin{enumerate}
\def\labelenumi{\arabic{enumi})}
\setcounter{enumi}{17}
\item
  Assume that G is compact. Let \(\varrho\) and \(\varrho'\) be unitary representations of G. The representation \(\varrho\) is weakly contained in \(\varrho'\) (cf. exercise 11 of V, p. 497) if and only if for every irreducible unitary representation \(\pi\) of G, the condition \(m_{\pi}(\varrho) \geqslant 1\) implies that \(m_{\pi}(\varrho') \geqslant 1\) . (Reduce to the case where \(\varrho'\) is irreducible; for every normalized diagonal matrix coefficient \(f\) of \(\varrho\) , prove by applying Exercise 13 of V, p.~498 that there exists a diagonal matrix coefficient of \(\varrho'\) which is not orthogonal to \(f\) , and conclude using the orthogonality relations.)
\item
  Let G be a compact group equipped with its normalized Haar measure \(\mu\) and let \(f \in \mathcal{C}(\mathrm{G})\) a continuous function on G. There exists an irreducible representation \(\pi \in \widehat{G}\) such that \(f = \dim(\pi)^{-1}\chi_{\pi}\) if and only if
\end{enumerate}

\[
\int_ {\mathrm{G}} f (x g x ^ {- 1} h) d \mu (x) = f (g) f (h)
\]

for all \((g,h)\in \mathrm{G}^2\)

\begin{enumerate}
\def\labelenumi{\arabic{enumi})}
\setcounter{enumi}{19}
\tightlist
\item
  Let H be a locally compact topological group equipped with a left Haar measure \(\nu\) . Let G be a compact subgroup of H. Let \(\varrho\) be an irreducible unitary representation of H in a Hilbert space E.
\end{enumerate}

We assume that the multiplicity of \(\pi\) in the restriction of \(\varrho\) to G is finite for every irreducible representation \(\pi \in \widehat{G}\) . We identify a function on G with a compactly supported measure on H.

\begin{enumerate}
\def\labelenumi{\alph{enumi})}
\item
  For every \(\pi \in \widehat{\mathrm{G}}\) , the endomorphism \(\varrho(f * \chi_{\pi})\) of E is of finite rank.
\item
  For every \(f \in \mathrm{L}^{1}(\mathrm{H})\) , the endomorphism \(\varrho(f)\) is compact.
\end{enumerate}

\begin{enumerate}
\def\labelenumi{\arabic{enumi})}
\setcounter{enumi}{20}
\tightlist
\item
  Let \(n \geqslant 1\) be an integer.
\end{enumerate}

\begin{enumerate}
\def\labelenumi{\alph{enumi})}
\tightlist
\item
  For every integer \(r > n^{2}\) and every family \((a_{1}, \ldots, a_{r})\) of elements of \(\mathcal{L}(\mathbf{C}^{n})\) , we have
\end{enumerate}

\[
\sum_ {\sigma \in \mathfrak {S} _ {r}} \varepsilon (\sigma) a _ {\sigma (1)} \dots a _ {\sigma (r)} = 0.\tag{4}
\]

\begin{enumerate}
\def\labelenumi{\alph{enumi})}
\setcounter{enumi}{1}
\tightlist
\item
  We denote by \(r(n)\) the smallest integer \(r \geqslant 1\) such that formula (4) is valid for every family \((a_i)_{1 \leqslant i \leqslant r}\) . We have \(r(n) \geqslant r(n-1) + 2\) (Consider a family \((a_1, \ldots, a_{r(n-1)-1})\) in \(\mathcal{L}(\mathbf{C}^{n-1})\) such that the expression above, denoted b, is nonzero; interpreting the elements of \(\mathcal{L}(\mathbf{C}^n)\) as matrices, choose i and j such that \(b_{i,j}\) is nonzero; view \(a_i\) and b as elements of \(\mathcal{L}(\mathbf{C}^n)\) by adding a row and a column of zeros, then define matrices \(a_{r(n-1)}\) and \(a_{r(n-1)+1}\) so that expression (4) for \((a_1, \ldots, a_{r(n-1)+1})\) is nonzero.)
\end{enumerate}

\begin{enumerate}
\def\labelenumi{\arabic{enumi})}
\setcounter{enumi}{21}
\tightlist
\item
  Let \(r \geqslant 1\) be an integer and let \(\mathrm{H} \subset \mathbf{GL}(r, \mathbf{C})\) be a closed subgroup.
\end{enumerate}

\begin{enumerate}
\def\labelenumi{\alph{enumi})}
\item
  The set of functions on H of the form \(g \mapsto \langle \varrho(g)x, \lambda \rangle\) , where \(\varrho\) is a continuous linear representation of H in a finite-dimensional complex vector space E, \(x \in E\) and \(\lambda \in E'\) , is a dense involutive subalgebra of \(\mathcal{C}(\mathrm{H})\) .
\item
  For any function \(f \in \mathcal{K}(\mathrm{H})\) nonzero, there exists a continuous linear representation \(\varrho\) of H in a finite-dimensional complex vector space such that \(\varrho(f)\) is nonzero.
\item
  Assume that H is a connected semisimple real Lie group. For every \(f \in \mathcal{K}(H)\) nonzero, there exists an irreducible continuous linear representation \(\varrho\) of H in a finite-dimensional complex vector space such that \(\varrho(f)\) is nonzero.
\end{enumerate}

\begin{enumerate}
\def\labelenumi{\arabic{enumi})}
\setcounter{enumi}{22}
\tightlist
\item
  We retain the hypotheses and notation of the preceding exercise, and assume that H is a connected semisimple real Lie group.
\end{enumerate}

Let G be a compact subgroup of H. Assume that there exists a connected solvable closed subgroup N of H such that H = GN. (For every semisimple real Lie group, this condition is satisfied when G is a maximal compact subgroup of H; see for example INT, VII, p.~91, § 3, n° 3, prop. 7 when G = SL(n, R).)

Let \(\pi \in \widehat{\mathbf{G}}\) and \(e_{\pi} = \dim (\pi)\overline{\chi}_{\pi}\) . Let \(\varrho\) be the irreducible representation of H given by the last question of the preceding exercise; we denote by E the space of \(\varrho\) .

\begin{enumerate}
\def\labelenumi{\alph{enumi})}
\item
  There exists a cyclic vector for the restriction of \(\varrho\) to G. (Use Lie's theorem, LIE, III, p.~241, cor., to find a nonzero \(x \in E\) such that C \(x\) is invariant under the action of N.) The multiplicity of \(\pi\) in the restriction of \(\varrho\) to G is at most equal to \(\dim(\pi)\) .
\item
  Let \(r\) be the integer \(r(\dim (\pi)^2)\) defined in Exercise 22; the identity
\end{enumerate}

\[
\sum_ {\sigma \in \mathfrak {S} _ {r}} \varepsilon (\sigma) a _ {\sigma (1)} \dots a _ {\sigma (r)} = 0
\]

is valid for every family \(a_{i}=e_{\pi}*g_{i}*e_{\pi}\) of elements of the algebra \(\mathcal{M}_{c}(\mathrm{H})\) of measures with compact support on H, where \(g_{i}\in\mathcal{K}(\mathrm{H})\) for \(1\leqslant i\leqslant r\) .

\begin{enumerate}
\def\labelenumi{\alph{enumi})}
\setcounter{enumi}{2}
\tightlist
\item
  Let \(\varpi\) be an irreducible unitary representation of H in a Hilbert space F. Let \(F_{\pi}\) be the \(\pi\)-isotypic component of the restriction of \(\varpi\) to G and let \(p\) be the orthogonal projector of F with image \(F_{\pi}\) . For every family \((u_i)_{1 \leqslant i \leqslant r}\) in \(\mathcal{L}(F_{\pi})\) , we have
\end{enumerate}

\[
\sum_ {\sigma \in \mathfrak {S} _ {r}} \varepsilon (\sigma) u _ {\sigma (1)} \dots u _ {\sigma (r)} = 0
\]

  (use Exercise 1 of V, p.~486).

\begin{enumerate}
\def\labelenumi{\alph{enumi})}
\setcounter{enumi}{3}
\tightlist
\item
  The multiplicity of \(\pi\) in the restriction of \(\varpi\) to G is at most equal to \(\dim(\pi)\) . (Use the result of Exercise 22 of V, p.~510.)
\end{enumerate}

\begin{enumerate}
\def\labelenumi{\arabic{enumi})}
\setcounter{enumi}{23}
\tightlist
\item
  A topological group G is said to be minimally almost periodic if every continuous homomorphism from G into a compact group is trivial.
\end{enumerate}

\begin{enumerate}
\def\labelenumi{\alph{enumi})}
\item
  A topological group G is minimally almost periodic if and only if G admits no non-trivial finite-dimensional unitary representation. If G is commutative, then G is minimally almost periodic if and only if every continuous unitary character of G is trivial.
\item
  Let G be a topological group and \(g \in G\) such that, for every integer \(k \geqslant 1\) , there exists an integer \(m \geqslant 1\) such that \(g^{km}\) is conjugate to g. Then one has \(\varrho(g) = 1_{\mathrm{E}}\) for every unitary representation \(\varrho\) of finite dimension of G in a Hilbert space E.
\item
  Let k be a subfield of C. The group \(\mathbf{SL}(2,k)\) endowed with the discrete topology is minimally almost periodic. (Verify that the preceding question applies to \(\begin{pmatrix}1 & z \\ 0 & 1\end{pmatrix}\) for all \(z \in k\) .)
\end{enumerate}

\begin{enumerate}
\def\labelenumi{\arabic{enumi})}
\setcounter{enumi}{24}
\tightlist
\item
  Let G be a locally compact group endowed with a left Haar measure \(\nu\) and \(K \subset G\) a compact subgroup of G endowed with the normalized Haar measure \(\mu\) .
\end{enumerate}

We denote by \(\mathcal{H}(\mathrm{G},\mathrm{K})\) the space of functions \(\varphi \in \mathcal{H}(\mathrm{G})\) such that \(\varphi (k_1gk_2) = \varphi (g)\) for all \(g\in \mathrm{G}\) and every \((k_{1},k_{2})\in \mathrm{K}\times \mathrm{K}\) .

The space \(\mathcal{H}(\mathrm{G},\mathrm{K})\) is a unital subalgebra of the algebra \(\mathcal{M}_c(\mathrm{G})\) ; one says that \((\mathrm{G},\mathrm{K})\) is a Gelfand pair if the algebra \(\mathcal{H}(\mathrm{G},\mathrm{K})\) is commutative. \(a)\) The map \(u^{\mathrm{K}}\) which to \(f\) associates the function \(f^{\mathrm{K}}\) defined by

\[
f ^ {\mathrm{K}} (g) = \int_ {\mathrm{K} \times \mathrm{K}} f (k _ {1} g k _ {2}) d (\mu \otimes \mu) (k _ {1}, k _ {2})
\]

is a projector in \(\mathcal{K}(\mathrm{G})\) whose image is \(\mathcal{H}(\mathrm{G},\mathrm{K})\) .

\begin{enumerate}
\def\labelenumi{\alph{enumi})}
\setcounter{enumi}{1}
\tightlist
\item
  Let \(f\in \mathcal{K}(\mathrm{G})\) . We have
\end{enumerate}

\[
\int_ {\mathrm{G}} f (g) d \nu (g) = \int_ {\mathrm{G}} f ^ {\mathrm{K}} (g) d \nu (g)
\]

and \((f\varphi)^{\mathrm{K}} = f^{\mathrm{K}}\varphi\) for every function \(\varphi \in \mathcal{H}(\mathrm{G},\mathrm{K})\) .

\begin{enumerate}
\def\labelenumi{\alph{enumi})}
\setcounter{enumi}{2}
\item
  Let \(C \subset G\) be a compact subset. There exists a function \(\varphi \in \mathcal{H}(G, K)\) such that \(\varphi(x) = 1\) for all \(x \in C\) .
\item
  Let (G,K) be a Gelfand pair. The group G is unimodular. (Note that the modulus of G belongs to \(\mathcal{H}(\mathrm{G},\mathrm{K})\) and interpret the integral of \(\varphi \in \mathcal{H}(\mathrm{G},\mathrm{K})\) as the value at \(e\) of a suitably chosen convolution.)
\item
  Suppose that there exists a continuous involution \(\theta\colon\mathrm{G}\to\mathrm{G}\) such that \(\mathrm{K}\theta(x)\mathrm{K}=\mathrm{K}x^{-1}\mathrm{K}\) for all \(x\in\mathrm{G}\) . Then (G,K) is a Gelfand pair. (Verify that the map \(f\mapsto\nu(f\circ\theta)\) on \(\mathcal{K}(\mathrm{G})\) is a left Haar measure on G, then that it is equal to \(\underline{\nu}\) ; deduce that \((f*g)\circ\theta=(f\circ\theta)*(g\circ\theta)\) and compare with the formula \(f*g=\check{g}*f\) .)
\item
  If G is compact, then \((\mathbf{G} \times \mathbf{G}, \Delta)\) is a Gelfand pair, where \(\Delta\) is the diagonal subgroup of \(G \times G\) image of the homomorphism \(x \mapsto (x, x)\) from G into \(G \times G\) .
\end{enumerate}

\begin{enumerate}
\def\labelenumi{\arabic{enumi})}
\setcounter{enumi}{25}
\tightlist
\item
  We resume the notation of the preceding exercise. We denote by \(\varphi_{\mathrm{K}}\) the characteristic function of K.
\end{enumerate}

\begin{enumerate}
\def\labelenumi{\alph{enumi})}
\item
  Let \(\pi\) be a unitary representation of G in a complex Hilbert space E. The representation of \(\mathcal{M}_c(\mathrm{G})\) in E defines, by passage to subspaces, a representation of the algebra \(\mathcal{H}(\mathrm{G},\mathrm{K})\) in the space \(\mathrm{E}^{\mathrm{K}}\) of K-invariant vectors of E.
\item
  If E admits a cyclic vector x such that \(x \in E^{K}\) , then \(\pi\) is irreducible.
\item
  Suppose that \(\pi\) is irreducible. If the space \(E^{K}\) is not reduced to 0, the representation of \(\mathcal{H}(G,K)\) in \(E^{K}\) is irreducible.
\end{enumerate}

\begin{enumerate}
\def\labelenumi{\alph{enumi})}
\setcounter{enumi}{3}
\tightlist
\item
  Let \(\varphi\) be a positive-definite function on G and let \((\varrho, x)\) be a cyclic Hilbertian realization of \(\varphi\) . We have \(\varphi \in \mathcal{H}(\mathrm{G}, \mathrm{K})\) if and only if \(x \in \mathrm{E}^{\mathrm{K}}\) .
\end{enumerate}

\begin{enumerate}
\def\labelenumi{\arabic{enumi})}
\setcounter{enumi}{26}
\tightlist
\item
  We resume the notation of the preceding exercise. We assume that (G, K) is a Gelfand pair.
\end{enumerate}

We say that \(\varphi\in\mathcal{H}(\mathrm{G},\mathrm{K})\) is a K-spherical function if the map

\[
f \mapsto \int_ {\mathrm{G}} f \check {\varphi} d \nu
\]

is a non-zero character of the algebra \(\mathcal{H}(\mathrm{G},\mathrm{K})\) .

\begin{enumerate}
\def\labelenumi{\alph{enumi})}
\tightlist
\item
  Let \(\varphi \in \mathcal{H}(\mathrm{G},\mathrm{K})\) . Prove that the following conditions are equivalent:
\end{enumerate}

\begin{enumerate}
\def\labelenumi{(\roman{enumi})}
\item
  The function \(\varphi\) is a K-spherical function;
\item
  For all \(x\) and \(y\) in \(\mathrm{G}\) , we have
\end{enumerate}

\[
\int_ {\mathrm{K}} \varphi (x k y) d \mu (k) = \varphi (x) \varphi (y);
\]

\begin{enumerate}
\def\labelenumi{(\roman{enumi})}
\setcounter{enumi}{2}
\tightlist
\item
  One has \(\varphi(e)=1\) and for every function \(\psi\in\mathcal{H}(\mathrm{G},\mathrm{K})\) , there exists \(c\in\mathbf{C}\) such that \(\psi*\varphi=c\varphi\) .
\end{enumerate}

\begin{enumerate}
\def\labelenumi{\alph{enumi})}
\setcounter{enumi}{1}
\item
  We denote by \(\mathcal{H}^1 (\mathrm{G},\mathrm{K})\) the space of \(f\in \mathrm{L}^1 (\mathrm{G})\) such that \(\varrho_{\mathrm{G}}(k_1,k_2)f = f\) for all \((k_{1},k_{2})\in \mathrm{K}\times \mathrm{K}\) . This is a Banach subalgebra of the involutive Banach algebra \(\mathrm{L}^1 (\mathrm{G})\) ; it is commutative.
\item
  The algebra \(\mathcal{H}(\mathrm{G},\mathrm{K})\) is identified with a dense subalgebra of \(\mathcal{H}^1 (\mathrm{G},\mathrm{K})\) .\\
\item
  The map which to \(\varphi \in \mathcal{H}(\mathrm{G},\mathrm{K})\) associates the linear form
\end{enumerate}

\[
f \mapsto \int_ {\mathrm{G}} f \check {\varphi} d \nu
\]

on \(\mathcal{H}^{1}(\mathrm{G},\mathrm{K})\) defines a bijection from the space of bounded K-spherical functions on G onto the space of non-zero characters of \(\mathcal{H}^{1}(\mathrm{G},\mathrm{K})\) . (Note that every character of \(\mathcal{H}^{1}(\mathrm{G},\mathrm{K})\) is continuous and of norm \(\leqslant 1\) , cf.~prop. 2 of I, p.~104.)

\begin{enumerate}
\def\labelenumi{\arabic{enumi})}
\setcounter{enumi}{27}
\tightlist
\item
  Let G be a locally compact group and K a compact subgroup of G.
\end{enumerate}

\begin{enumerate}
\def\labelenumi{\alph{enumi})}
\item
  The pair (G, K) is a Gelfand pair if and only if, for every unitary representation \(\pi\) of G in a complex Hilbert space E, the dimension of \(E^{K}\) is at most 1. (To prove that the condition is necessary, prove that an irreducible representation of a commutative involutive Banach algebra in a Hilbert space is of dimension 1; to prove that it is sufficient, use the cor. of th. 4 of V, p.~454 to verify that the characters of the algebra \(\mathcal{H}(G,K)\) separate its points.)
\item
  We assume that (G, K) is a Gelfand pair. Let \(\varphi \in \mathcal{H}(\mathrm{G},\mathrm{K})\) . Assume that \(\varphi\) is a positive-definite function and that \(\varphi(e) = 1\) . Show that \(\varphi\) is a K-spherical function if and only if \(\varphi\) has a Hilbert representation \((\varrho, x)\) such that \(\varrho\) is irreducible.
\item
  Suppose that \((\mathrm{G},\mathrm{K})\) is a Gelfand pair. The set of positive-definite K-spherical functions on G is in bijection with the set of irreducible representations \(\pi\in\widehat{G}\) such that \(E_{\pi}^{K}\) is nonzero.
\end{enumerate}

\begin{enumerate}
\def\labelenumi{\arabic{enumi})}
\setcounter{enumi}{28}
\tightlist
\item
  Let G be a compact group and K a closed subgroup of G. Suppose that (G, K) is a Gelfand pair. We denote by \(\nu\) (resp. \(\mu\) ) the normalized Haar measure on G (resp. on K).
\end{enumerate}

Let \(\varphi\) be a K-spherical function on G.

\begin{enumerate}
\def\labelenumi{\alph{enumi})}
\item
  We denote by \(E_{\varphi}\) the subspace of \(\mathcal{C}(G)\) functions of the form \(\alpha * \varphi\) , where \(\alpha\) is a measure on G with finite support. For every \(\psi \in \mathcal{H}(G, K)\) , the endomorphism \(v_{\psi}: f \mapsto f * \psi\) of \(\mathcal{C}(G)\) is compact and \(E_{\varphi}\) is contained in an eigenspace of \(v_{\psi}\) .
\item
  The space \(E_{\varphi}\) is of finite dimension. It identifies with a subspace of G-finite vectors of the regular representation \(\gamma_{G}\) containing \(\varphi\) .
\item
  The space of K-invariant vectors of \(E_{\varphi}\) is of dimension 1 and the representation of G in \(E_{\varphi}\) is irreducible.
\end{enumerate}

\begin{enumerate}
\def\labelenumi{\alph{enumi})}
\setcounter{enumi}{3}
\tightlist
\item
  Let \((f_i)_{i\in I}\) be an orthonormal basis of \(\mathrm{E}_{\varphi}\) and \(k\colon \mathbf{G}\times \mathbf{G}\to \mathbf{C}\) the function defined by

\[
k (x, y) = \sum_ {i \in \mathrm{I}} \overline {{f _ {i} (x)}} f _ {i} (y).
\]

One has

\[
f (x) = \int_ {\mathrm{G}} k (x, y) f (y) d \nu (y)
\]

for all \(f \in \mathrm{E}_{\varphi}\) and \(x \in \mathrm{G}\) .

\end{enumerate}

\begin{enumerate}
\def\labelenumi{\alph{enumi})}
\setcounter{enumi}{4}
\item
  One has \(k(x,y)=\dim(\mathrm{E}_{\varphi})\varphi(y^{-1}x)\) for all \((x,y)\in\mathrm{G}\times\mathrm{G}\) .
\item
  The function \(\varphi\) is of positive type.
\item
  Let \(\varphi\) be a K-spherical function on G and let \((\pi, x)\) be a Hilbert realization of \(\varphi\) . The representation \(\pi\) is irreducible; we have
\end{enumerate}

\[
\varphi (x) = \int_ {\mathrm{K}} \chi_ {\pi} (x ^ {- 1} k) d \mu (k)
\]

for all \(x \in \mathrm{G}\) , and \(\int_{\mathrm{G}} |\varphi|^2 d\nu = \dim(\pi)^{-1}\) .

\begin{enumerate}
\def\labelenumi{\arabic{enumi})}
\setcounter{enumi}{29}
\tightlist
\item
  Let \(n \in N\) . Let us denote by G the subgroup of the affine group \(\mathrm{A}(\mathbf{R}^{n})\) of the space \(R^{n}\) formed by the elements \(g \in \mathrm{A}(\mathbf{R}^{n})\) such that the linear map \(g_{0}\) associated with g belongs to \(\mathbf{SO}(\mathbf{R}^{n})\) ; it is a real Lie group (cf. A, II, p. 131 and LIE, III, p. 102, example). We denote by \(g = (g_{0}, a)\) the elements of G, with \(g_{0} \in \mathbf{SO}(\mathbf{R}^{n})\) and \(a \in R^{n}\) .
\end{enumerate}

\begin{enumerate}
\def\labelenumi{\alph{enumi})}
\item
  The pair \((\mathbf{G},\mathbf{SO}(\mathbf{R}^n))\) is a Gelfand pair.
\item
  Let \(\varphi \in \mathcal{H}(\mathrm{G},\mathbf{SO}(\mathbf{R}^n))\) . Denote by \(f\colon \mathbf{R}^n\to \mathbf{C}\) the function defined by \(f(r) = \varphi (1_{\mathbf{R}^N},g(e_1))\) . The map \(u\colon \varphi \mapsto f\) is an injective morphism of complex algebras from \(\mathcal{H}(\mathrm{G},\mathrm{K})\) into \(\mathcal{K}(\mathbf{R}^n)\) ; its image is the subspace of \(\mathcal{K}(\mathbf{R}^n)\) consisting of functions \(f\) such that \(f(g_0(x)) = f(x)\) for all \(g_0\in \mathbf{SO}(\mathbf{R}^n)\) and every \(x\in \mathbf{R}^n\) .
\item
  We denote by \(\Delta\) the Laplacian on \(\mathbf{R}^n\) . A function \(\varphi \in \mathcal{K}(\mathrm{G})\) is a spherical function if and only if the function \(f = u(\varphi)\) is of class \(\mathrm{C}^{\infty}\) , satisfies \(f(0) = 1\) , and if there exists \(\lambda \in \mathbf{C}\) such that \(\Delta f = \lambda f\) (Use the fact that \(\Delta(f \circ g_0) = (\Delta f) \circ g_0\) for all \(g_0 \in \mathbf{SO}(\mathbf{R}^n)\) .)
\end{enumerate}

\begin{enumerate}
\def\labelenumi{\arabic{enumi})}
\setcounter{enumi}{30}
\tightlist
\item
  Let \(n\) be an integer \(\geqslant 1\) . We denote by N the Euclidean norm on \(\mathbf{R}^n\) , and we denote by \(\lambda_n\) the Lebesgue measure on \(\mathbf{R}^n\) . We denote by \(\omega_n\) the unique positive measure of total mass 1 on the sphere \(\mathbf{S}_n \subset \mathbf{R}^{n+1}\) invariant under the action of the orthogonal group \(\mathbf{O}(n+1, \mathbf{R})\) (cf.~INT, VIII, p.~116, exercise 8).
\end{enumerate}

We identify \(R^{n}\) with the complement of a projective hyperplane H in the projective space \(\mathbf{P}_{n}(\mathbf{R})\) (cf.~TG, VI, p.~15, prop. 5), and we shall interpret the measures on \(R^{n}\) as measures on \(\mathbf{P}_{n}(\mathbf{R})\) supported on the complement of H.

For every \(g \in \mathbf{GL}(n+1, \mathbf{R})\) , we denote by \(\widetilde{g}\) the homeomorphism of \(\mathbf{P}_{n}(\mathbf{R})\) defined by g.

\begin{enumerate}
\def\labelenumi{\alph{enumi})}
\tightlist
\item
  The positive measure \(c_{n}\) on \(\mathbf{R}^n\) defined by
\end{enumerate}

\[
c _ {n} = \pi^ {- (n + 1) / 2} \Gamma \Big (\frac {n + 1}{2} \Big) (1 + \mathrm{N} ^ {2}) ^ {- (n + 1) / 2} \lambda_ {n}
\]

(“Cauchy measure”) is bounded and has total mass 1. If \(n = 1\) this is the stable measure of parameter 1 (cf.~exercise V, p.~496). For every projective hyperplane \(\mathrm{H}'\) of \(\mathbf{P}_n(\mathbf{R})\) , we have \(c_n(\mathrm{H}') = 0\) .

\begin{enumerate}
\def\labelenumi{\alph{enumi})}
\setcounter{enumi}{1}
\item
  The canonical map from \(\mathbf{R}^{n+1} - \{0\}\) into \(\mathbf{P}_n(\mathbf{R})\) defines, by passing to subspaces, a continuous surjective map \(\sigma: \mathbf{S}_n - (\mathbf{R}^n \times \{0\}) \to \mathbf{R}^n\) . We have \(\sigma(\omega_n) = c_n\) .
\item
  Let \(g \in \mathbf{GL}(n+1, \mathbf{R})\) . There exists \(a \in \mathbf{GL}(n, \mathbf{R})\) and \(b \in R^n\) such that the image measure \(\widetilde{g}(c_n)\) is equal to the image of \(c_n\) by the affine map \(x \mapsto a(x) + b\) .
\end{enumerate}

In what follows, we propose to prove a converse of this assertion.

\begin{enumerate}
\def\labelenumi{\alph{enumi})}
\setcounter{enumi}{3}
\tightlist
\item
  Let \(\mu\) be a bounded positive measure on \(\mathbf{P}_{n}(\mathbf{R})\) such that every projective hyperplane is \(\mu\) -negligible. The set of \(g \in \mathbf{GL}(n+1, \mathbf{R})\) such that \(g(\mu) = \mu\) and \(|\det(g)| = 1\) is a compact subgroup of \(\mathbf{GL}(n+1, \mathbf{R})\) .

\end{enumerate}

\begin{enumerate}
\def\labelenumi{\alph{enumi})}
\setcounter{enumi}{4}
\item
  Let \(\mu\) be a bounded positive measure on \(\mathbf{P}_n(\mathbf{R})\) of total mass 1 such that, for all \(g \in \mathbf{GL}(n + 1, \mathbf{R})\) there exists an affine map \(f\) of \(\mathbf{R}^n\) such that \(g(\mu) = f(\mu)\) . Every projective hyperplane is \(\mu\) -negligible.
\item
  Let \(K_{\mu}\) be the subgroup of \(g \in \mathbf{GL}(n+1, \mathbf{R})\) such that \(g(\mu) = \mu\) . There exists \(h \in \mathbf{GL}(n+1, \mathbf{R})\) such that \(K_{\mu} \subset h \mathbf{O}(n+1, \mathbf{R}) h^{-1}\) .
\item
  Let \(K = h^{-1}K_{\mu}h\) . We have \(\mathbf{O}(n + 1, \mathbf{R}) = \mathbf{G} \cdot \mathbf{K}\) , where G is the subgroup of \(\mathbf{O}(n + 1, \mathbf{R})\) consisting of the matrices
\end{enumerate}

\[
\left( \begin{array}{c c} x & 0 \\ 0 & 1 \end{array} \right)
\]

with \(x\in \mathbf{O}(n,\mathbf{R})\)

\begin{enumerate}
\def\labelenumi{\alph{enumi})}
\setcounter{enumi}{7}
\tightlist
\item
  The measure \(h^{-1}(\mu)\) coincides with \(c_{n}\) . (Observe that the group K acts transitively on \(S_{n}\) and that the image measure \(\sigma(h^{-1}(\mu))\) is a K-invariant measure on \(S_{n}\) .)
\end{enumerate}

\begin{enumerate}
\def\labelenumi{\arabic{enumi})}
\setcounter{enumi}{31}
\tightlist
\item
  Let G be a locally compact group. Suppose there exists a subgroup H of G that is compact and open in G. We denote by \(\mu\) the unique left Haar measure on G such that \(\mu(\mathrm{H}) = 1\) and \(\varphi\) the characteristic function of H.
\end{enumerate}

\begin{enumerate}
\def\labelenumi{\alph{enumi})}
\item
  Let \(\pi\) be a square-integrable representation of G in a complex Hilbert space E. Show that the endomorphism \(\pi(\varphi)\) is a Hilbert-Schmidt endomorphism.
\item
  The space \(E^{H}\) is finite-dimensional and satisfies
\end{enumerate}

\[
\dim (\mathrm{E} ^ {\mathrm{H}}) \leqslant \frac {1}{c (\pi)}
\]

where \(c(\pi)\) is the formal degree of \(\pi\) relative to \(\{e\}\) (Use Exercise 6 of V, p.~487.)

\begin{enumerate}
\def\labelenumi{\alph{enumi})}
\setcounter{enumi}{2}
\tightlist
\item
  A discrete infinite group has no square-integrable representation; if \(n \geqslant 2\) is an integer, the group \(\mathbf{SL}(n, \mathbf{Z})\) has no square-integrable representation modulo the center.
\end{enumerate}

